\chapter{Initial Implementation}

Trong giai đoạn triển khai ban đầu, chúng tôi tập trung vào việc thiết lập nền tảng cốt lõi cho hệ thống Multi-Agent, bao gồm việc định nghĩa cấu trúc trạng thái (State Schema), chuẩn bị các công cụ tích hợp (Tool Binding), và cấu hình môi trường chạy cục bộ với Ollama và LangGraph.

\begin{figure}[H]
\vspace{2pt}
\centering
% Bạn có thể dùng lại hình sơ đồ kiến trúc LangGraph ở chương trước hoặc vẽ mới
\includegraphics[scale=0.5]{Images/Implement/agent_workflow.png} 
\vspace{5mm}
\caption{Luồng dữ liệu trong hệ thống Multi-Agent}
\end{figure}

\section{Data Structure Design (Thiết kế Cấu trúc Dữ liệu)}

Khác với các hệ thống RAG truyền thống dựa vào Cơ sở dữ liệu Vector, hệ thống Agentic của chúng em hoạt động dựa trên luồng trạng thái (Stateful Workflow). Do đó, việc thiết kế cấu trúc dữ liệu tập trung vào đối tượng **Shared State** - bộ nhớ chia sẻ giữa các tác nhân.

\subsection{Agent State Schema}

Dựa trên thư viện \texttt{LangGraph}, chúng em định nghĩa một lược đồ trạng thái toàn cục (Global State Schema) để duy trì ngữ cảnh xuyên suốt phiên kiểm toán. Cấu trúc này được cài đặt dưới dạng Python \texttt{TypedDict}:

\begin{itemize}
    \item \textbf{Input Data:}
    \begin{itemize}
        \item \texttt{user\_request}: Chuỗi văn bản chứa yêu cầu ban đầu của người dùng (ví dụ: "Scan S3 services").
        \item \texttt{scan\_plan}: Đối tượng JSON chứa cấu hình phạm vi quét do \textit{Planning Agent} sinh ra.
    \end{itemize}
    
    \item \textbf{Process Data:}
    \begin{itemize}
        \item \texttt{job\_ids}: Danh sách các định danh tiến trình (Process IDs) của Prowler đang chạy ngầm.
        \item \texttt{raw\_findings}: Dữ liệu thô thu được từ công cụ quét, chưa qua xử lý.
    \end{itemize}
    
    \item \textbf{Output Intelligence:}
    \begin{itemize}
        \item \texttt{enriched\_findings}: Dữ liệu đã được \textit{Risk Agent} làm giàu thông tin và chấm điểm rủi ro.
        \item \texttt{failed\_findings}: Danh sách lọc gọn các lỗi nghiêm trọng cần khắc phục.
        \item \texttt{remediation\_results}: Kết quả trả về từ các hành động sửa lỗi tự động.
    \end{itemize}
\end{itemize}

\begin{table}[H]
\centering
\label{tab:state_schema}
\begin{tabular}{p{4cm}p{3cm}p{7cm}}
\hline
\textbf{Field Name} & \textbf{Data Type} & \textbf{Description} \\ \hline
\texttt{scan\_plan} & \texttt{Dict[str, Any]} & Cấu hình quét bao gồm danh sách dịch vụ và regions. \\ \hline
\texttt{enriched\_findings} & \texttt{List[Dict]} & Danh sách lỗi kèm theo \texttt{severity} và \texttt{risk\_score}. \\ \hline
\texttt{remediation\_needed} & \texttt{bool} & Cờ (Flag) điều hướng luồng xử lý: \texttt{True} nếu cần kích hoạt Remediation Agent. \\ \hline
\end{tabular}
\caption{Các trường dữ liệu chính trong AgentState.}
\end{table}

\section{Data Preparation \& Tooling (Chuẩn bị Công cụ)}

Thay vì thu thập dữ liệu tĩnh, hệ thống Agent cần chuẩn bị các "giác quan" và "tay chân" để tương tác với môi trường AWS. Quá trình này bao gồm việc đóng gói các công cụ CLI và SDK thành các hàm Python có thể gọi được (Callable Functions).

\textbf{Tool Definition Steps:}

\begin{itemize}
    \item \textbf{Prowler Integration:} 
    Chúng tôi xây dựng một wrapper xung quanh công cụ Prowler để Agent có thể kích hoạt quét theo dòng lệnh.
    \begin{lstlisting}[language=Python, caption=Prowler Tool Wrapper]
def run_prowler_scan(service: str, region: str = "us-east-1"):
    cmd = f"prowler aws --services {service} --region {region} --output-format json"
    \end{lstlisting}

    \item \textbf{Remediation Tools (Boto3):} 
    Các hành động sửa lỗi được định nghĩa sẵn bằng thư viện \texttt{boto3}. Mỗi tool được gán một mô tả (docstring) chi tiết để LLM hiểu được ngữ nghĩa và cách sử dụng.
    Ví dụ: \texttt{s3\_block\_public\_access}, \texttt{iam\_rotate\_access\_key}.
    
    \item \textbf{Tool Registry:} 
    Tất cả các công cụ được đăng ký vào một danh sách \texttt{SCANNER\_AGENT\_TOOLS} và \texttt{ALL\_TOOLS} (trong file \texttt{agent\_tools.py}) để phục vụ việc liên kết động (Dynamic Binding) với mô hình ngôn ngữ.
\end{itemize}

\section{Data Preprocessing (Tiền xử lý Dữ liệu Quét)}

Dữ liệu đầu ra từ Prowler thường rất lớn và chứa nhiều thông tin nhiễu kỹ thuật. Để LLM có thể xử lý hiệu quả, chúng tôi cài đặt quy trình tiền xử lý tại \textit{Monitoring Agent} và \textit{Risk Evaluation Agent}.

\textbf{Key Preprocessing Steps:}

\begin{itemize}
    \item \textbf{Normalization (Chuẩn hóa):} 
    Chuyển đổi các định dạng JSON khác nhau của Prowler (OCSF hoặc Native) về một cấu trúc phẳng (flat structure) thống nhất.
    
    \item \textbf{Filtering (Lọc nhiễu):} 
    Loại bỏ các findings có trạng thái \texttt{PASS} hoặc \texttt{INFO}, chỉ giữ lại các findings có trạng thái \texttt{FAIL}. Logic này được thực hiện bằng mã nguồn Python thuần túy để giảm tải cho LLM và tăng tốc độ xử lý.
    
    \item \textbf{Field Extraction (Trích xuất trường):} 
    Từ dữ liệu thô, chúng tôi chỉ trích xuất các trường quan trọng cần thiết cho việc ra quyết định:
    \begin{itemize}
        \item \texttt{check\_id}: Mã định danh lỗi (ví dụ: \texttt{s3\_bucket\_public\_access}).
        \item \texttt{resource\_arn}: Định danh tài nguyên bị ảnh hưởng.
        \item \texttt{status\_extended}: Mô tả chi tiết nguyên nhân lỗi.
    \end{itemize}
\end{itemize}

Đoạn mã dưới đây minh họa logic làm sạch dữ liệu trong \texttt{risk\_evaluation\_agent.py}:

\begin{lstlisting}[language=Python, caption=Finding Cleaning Logic]
def _clean_output_for_assessment(self, finding: dict, risk_data: dict) -> dict:
    return {
        "finding_id": finding.get("finding_id"),
        "service_group": service_group,
        "issue_hint": str(finding.get("prowler_check_id") or "").strip(),
        "status": finding.get("status"), 
        "severity": risk_data.get("severity"),
        "risk_score": risk_data.get("risk_score")
    }
\end{lstlisting}

\section{Orchestrator Setup (Thiết lập Bộ điều phối)}

Thay vì thiết lập Database Server, "cơ sở dữ liệu" của hệ thống chính là Graph Workflow. Chúng tôi sử dụng \texttt{StateGraph} của LangGraph để kết nối các Agent.

\subsection{Node \& Edge Definition}

Hệ thống được khởi tạo bằng cách đăng ký các hàm xử lý (Node) và định nghĩa luồng đi (Edge):

\begin{itemize}
    \item \textbf{Nodes:} Bao gồm \texttt{plan\_node}, \texttt{scan\_node}, \texttt{monitor\_node}, \texttt{risk\_eval\_node}, \texttt{remediate\_node}, v.v. Mỗi node là một hàm Python nhận vào \texttt{AgentState} và trả về phần cập nhật (State Update).
    \item \textbf{Conditional Edges:} Logic rẽ nhánh quan trọng được cài đặt tại \texttt{human\_approval\_node}.
    \[
    \text{Next Node} = \begin{cases} 
    \text{remediate} & \text{if } \text{approved} == \text{True} \\
    \text{report} & \text{if } \text{approved} == \text{False}
    \end{cases}
    \]
\end{itemize}

\subsection{File-based Persistence}
Để đảm bảo tính bền vững (Persistence) mà không cần DB phức tạp, hệ thống lưu các snapshot trạng thái xuống thư mục \texttt{data/} dưới dạng JSON (\texttt{pre\_scan.json}, \texttt{post\_scan.json}).

\section{Local LLMs Setup (Cấu hình Mô hình Ngôn ngữ)}

Chúng tôi sử dụng \textbf{Ollama} làm nền tảng chạy mô hình ngôn ngữ cục bộ, đảm bảo dữ liệu nhạy cảm về hạ tầng AWS không bị gửi ra ngoài Internet.

\subsection{Model Selection & Configuration}
Mô hình được lựa chọn là \texttt{llama3.2} (phiên bản 3B parameters) nhờ khả năng cân bằng giữa tốc độ và khả năng suy luận (Reasoning).

\begin{verbatim}
OLLAMA_MODEL = "llama3.2"
OLLAMA_BASE_URL = "http://localhost:11434"
\end{verbatim}

\subsection{Integration with Agents}

Việc tích hợp LLM vào từng Agent được thực hiện thông qua thư viện \texttt{langchain\_ollama}. Chúng tôi áp dụng hai chiến lược cấu hình khác nhau tùy thuộc vào vai trò của Agent:

\begin{enumerate}
    \item \textbf{JSON Mode (cho Risk Evaluation Agent):}
    Cấu hình \texttt{format="json"} và \texttt{temperature=0} để đảm bảo đầu ra luôn là cấu trúc dữ liệu máy có thể đọc được, phục vụ cho việc chấm điểm rủi ro.
    
    \item \textbf{Tool Calling Mode (cho Scanner \& Remediate Agent):}
    Sử dụng phương thức \texttt{bind\_tools} để cho phép mô hình nhận biết và lựa chọn công cụ từ danh sách \texttt{agent\_tools}.
    
\begin{lstlisting}[language=Python, caption=LLM Initialization with Tools]
self.llm = ChatOllama(model=model_name, temperature=0)
self.llm_with_tools = self.llm.bind_tools(SCANNER_AGENT_TOOLS)
\end{lstlisting}
\end{enumerate}

Cấu hình này cho phép hệ thống vận hành trơn tru trên máy cá nhân với RAM 8GB-16GB mà vẫn đảm bảo khả năng xử lý các tác vụ phức tạp.




